\documentclass[../DoAn.tex]{subfiles}
\begin{document}

\section{Kết luận}

Đồ án đã nghiên cứu bài toán lập lịch và định tuyến lấy hàng cho hệ thống vận tải phối hợp xe tải và drone có cửa sổ thời gian, trong đó nhiều xe tải và nhiều drone cùng hoạt động, mỗi phương tiện có thể thực hiện nhiều chuyến và mục tiêu là tối thiểu hóa makespan. Bài toán được phát biểu hình thức với đầy đủ các ràng buộc về loại phương tiện, tải trọng, thời gian bay của drone, cửa sổ thời gian và thời gian lưu hàng trên phương tiện. Đồ án cũng chứng minh bài toán là NP-khó bằng phép quy dẫn từ bài toán người du lịch, qua đó lý giải việc lựa chọn hướng tiếp cận metaheuristic.

Giải pháp đề xuất gồm hai giai đoạn. Giai đoạn thứ nhất xây dựng lời giải ban đầu bằng thuật toán Greedy Insertion theo thứ tự thời hạn sớm nhất. Giai đoạn thứ hai cải thiện lời giải bằng Adaptive Tabu Search với sáu toán tử lân cận, trong đó toán tử Trip-relocate khai thác trực tiếp cấu trúc chuỗi chuyến của từng phương tiện. Thuật toán dùng hàm mục tiêu chuẩn hóa có bốn thành phần vi phạm với hệ số phạt được điều chỉnh độc lập bằng Strategic Oscillation, danh sách Tabu theo ba loại thuộc tính, và cơ chế Ruin \& Recreate có định hướng vào các khách vi phạm và các phương tiện critical. Bên cạnh đó, đồ án cài đặt bốn chiến lược sinh lân cận và phân tích độ phức tạp của chúng, cho thấy các chiến lược chỉ chọn một toán tử mỗi vòng có chi phí mỗi vòng lặp tăng tuyến tính theo số khách.

Thực nghiệm trên 80 instance có từ 6 đến 100 khách cho thấy thuật toán luôn tìm được lời giải khả thi, và toàn bộ 800 lời giải đều được bộ kiểm tra độc lập xác nhận hợp lệ. Thuật toán đạt đúng makespan của lời giải tham chiếu ở 35 trên 36 instance có từ 6 đến 12 khách, với thời gian giải dưới một giây. Trên toàn bộ dữ liệu, lời giải tốt nhất qua năm lần chạy chênh lệch trung bình khoảng 1,9\% so với tham chiếu, và ở một instance còn tốt hơn tham chiếu. So sánh qua năm lần chạy độc lập cho thấy chiến lược chọn toán tử tuần tự cho chất lượng tốt hơn nhẹ so với chiến lược ngẫu nhiên, với tổng gap $-5{,}65\%$ theo lời giải tốt nhất và $-6{,}98\%$ theo lời giải trung bình trên 80 instance, trong khi chiến lược ngẫu nhiên nhanh hơn khoảng 9\%. Khác biệt giữa hai chiến lược chỉ xuất hiện ở các bài toán từ 20 khách trở lên và nhỏ hơn độ dao động giữa các lần chạy.

Bên cạnh các kết quả trên, đồ án vẫn còn một số hạn chế. Thứ nhất, chênh lệch so với lời giải tham chiếu tăng theo kích thước bài toán, tới khoảng 5\% với 100 khách, và có một số instance chênh lệch lớn trên 15\%. Thứ hai, lời giải tham chiếu được tính với một số thiết lập khác với bài toán của đồ án, nên chưa đánh giá được chính xác khoảng cách tới lời giải tối ưu. Thứ ba, số lần chạy cho mỗi chiến lược mới là năm lần, chưa đủ để kiểm định ý nghĩa thống kê của khác biệt giữa hai chiến lược. Thứ tư, chương trình vẫn sao chép và tính lại toàn bộ các phương tiện bị ảnh hưởng khi đánh giá mỗi move, chưa dùng kỹ thuật đánh giá tăng dần, nên chi phí đánh giá còn cao với các bài toán lớn.

\section{Hướng phát triển trong tương lai}

Từ các hạn chế trên, đồ án đề xuất một số hướng phát triển. Hướng thứ nhất là tăng tốc độ đánh giá move bằng kỹ thuật đánh giá tăng dần, chẳng hạn tính trước độ trễ cho phép (forward time slack) của từng vị trí trong chuyến theo cách của Savelsbergh~\cite{savelsbergh1985local}, để kiểm tra ràng buộc cửa sổ thời gian trong thời gian hằng số. Khi đó có thể tăng số move được đánh giá mỗi vòng hoặc số vòng lặp trong cùng ngân sách thời gian. Hướng thứ hai là đánh giá đầy đủ chiến lược chọn toán tử thích nghi và chiến lược duyệt toàn bộ sáu toán tử với cùng ngân sách thời gian, đồng thời tăng số lần chạy và áp dụng các kiểm định thống kê như kiểm định Wilcoxon để so sánh các chiến lược một cách chặt chẽ. Hướng thứ ba là xây dựng mô hình quy hoạch nguyên hỗn hợp cho bài toán và giải chính xác các instance nhỏ, hoặc tính cận dưới cho các instance lớn, để có mốc đánh giá chính xác hơn. Hướng thứ tư là song song hóa việc đánh giá tập ứng viên trên nhiều luồng, vì các ứng viên trong cùng một vòng lặp được đánh giá độc lập với nhau. Cuối cùng, mô hình có thể được mở rộng để phù hợp hơn với thực tế, chẳng hạn xét thời gian phục vụ tại khách, thời gian thay pin của drone giữa các chuyến, hoặc yêu cầu lấy hàng xuất hiện động trong ca làm việc.

\end{document}
